\section{Reductions for the exact four-regular bound}\label{sec:four-regular}

All graphs in this section are loopless multigraphs.
Parallel edges are distinct, and two parallel edges form a cycle.
Write $c(G)$ for the number of cycles.
The target is
\begin{equation}\label{eq:four-target}
 c(G)\le 2^n+n
 \qquad\text{for every four-regular graph on $n$ vertices.}
\end{equation}
This is Open Problem~2 of Alahmadi, Aldred, dela Cruz, Sol\'e, and Thomassen,
\emph{The maximum number of minimal codewords in long codes},
\url{https://arxiv.org/html/1201.6043}, Section~6.
A focused literature search found no subsequent resolution.
This search does not establish that none exists.
The results below do not settle \eqref{eq:four-target}.

\subsection{Complementation and the possible counterexamples}

For connected $G$, its cycle space has $2^{n+1}$ elements.
The pairs $\{X,E(G)\setminus X\}$ number $2^n$.
Let $a$ count pairs containing two cycles, and let $b$ count pairs containing no cycle.
Then
\[
 c(G)=2^n+a-b.
\]
Both members of a pair counted by $a$ must be Hamilton cycles.
Indeed, their degrees sum to four at every vertex.
Each cycle contributes at most two there, so both must contribute two.
Consequently, a graph without a Hamilton decomposition has $c(G)\le2^n$.
This is the specialization of the cited paper's Theorem~6.

A graph with a Hamilton decomposition is four-edge-connected.
Each Hamilton cycle crosses every nontrivial cut at least twice.
For $n\ge3$, such a graph has edge multiplicity at most two.
Each Hamilton cycle uses at most one edge between a given pair of vertices.

\subsection{Doubled-edge expansion preserves the proposed bound}

Suppose that $G$ has exactly two edges between $u$ and $v$.
Replace these edges by two $uw$ edges and two $wv$ edges, where $w$ is new.
Call the resulting graph $G'$.
This operation preserves four-regularity and adds one vertex.
Let $a_{uv}(G)$ count cycles using exactly one of the two $uv$ edges.

\begin{proposition}\label{prop:double-expansion}
For connected four-regular $G$ on $n\ge3$ vertices,
\[
 c(G')=c(G)+a_{uv}(G)+1,
 \qquad a_{uv}(G)\le2^n.
\]
Thus \eqref{eq:four-target} for $G$ implies \eqref{eq:four-target} for $G'$.
Moreover, $a_{uv}(G)=2^n$ holds if and only if $G$ is the doubled $n$-cycle.
\end{proposition}

\begin{proof}
Cycles avoiding both $uv$ edges remain unchanged.
Each cycle using exactly one gains twice as many choices after expansion.
The original two-edge cycle is replaced by two two-edge cycles.
This proves the identity.

Let $N$ result from deleting the two $uv$ edges.
If $N$ is disconnected, its terminals lie in different components, and $a_{uv}(G)=0$.
Otherwise $\beta(N)=n-1$, and
\[
 a_{uv}(G)=2P_{uv}(N)\le2\cdot2^{\beta(N)}=2^n.
\]
Here the path bound follows by injecting paths into the affine space of $\{u,v\}$-joins.
The argument allows parallel edges.

Suppose equality holds.
Every join of $N$ is then a path.
For any such path $A$, deleting $E(A)$ leaves a forest.
Otherwise $A$ together with an edge-disjoint cycle would be a join that is not a path.
Thus $|E(A)|\ge\beta(N)=n-1$, and every terminal path is Hamiltonian.
Fix one such path.
An edge joining nonconsecutive vertices of that path would give a shorter terminal path.
Therefore all edges of $N$ join consecutive vertices of the fixed path.
The terminal degrees are two, and every internal degree is four.
Successive degree equations force every consecutive pair to have multiplicity two.
Adding the deleted pair makes $G$ a doubled cycle.
The converse follows by counting the $2^{n-1}$ terminal paths in a doubled path.
\end{proof}

More generally, replace the doubled edge by a doubled path with $t$ new internal vertices.
Call this graph $G_t$, with $G_0=G$.
Let $b_{uv}(G)$ count cycles avoiding both $uv$ edges.
Direct counting gives
\begin{align}
 c(G_t)&=2^t a_{uv}(G)+b_{uv}(G)+t+1,\\
 \bigl(2^{n+t}+n+t-c(G_t)\bigr)
 &=\bigl(2^n+n-c(G)\bigr)+(2^t-1)\bigl(2^n-a_{uv}(G)\bigr).
 \label{eq:double-chain-deficit}
\end{align}
If $G$ is not a doubled cycle, then $a_{uv}(G)$ is even and at most $2^n-2$.
Hence each expansion sequence increases the deficit by at least $2(2^t-1)$.
This conclusion assumes only the displayed initial deficit, not the full conjecture.

\begin{corollary}\label{cor:four-core}
If \eqref{eq:four-target} fails, a counterexample of minimum order has all the following properties:
it is connected, Hamilton-decomposable, and four-edge-connected.
Its underlying simple graph has minimum degree at least three.
Its parallel pairs form a matching.
\end{corollary}

\begin{proof}
Disconnected graphs cannot be minimum counterexamples.
Apply the smaller-order bounds to their components and use $\sum_i2^{n_i}\le2^{\sum_i n_i}$.
Complementation gives the Hamilton decomposition, which gives four-edge-connectivity.

Suppose a vertex $w$ has only two distinct neighbors, $u$ and $v$.
Each Hamilton cycle uses one edge to each neighbor at $w$.
Thus both incident multiplicities are two.
For $n>3$, there is no $uv$ edge.
Such an edge would close a triangle in whichever Hamilton cycle contains it.
Suppress $w$ and replace its four edges by a doubled $uv$ edge.
The resulting graph is loopless, four-regular, connected, and has one fewer vertex.
Proposition~\ref{prop:double-expansion} then excludes a minimum counterexample.
The smaller base cases are the doubled triangle and the four parallel edges on two vertices.
Their counts are $11$ and $6$.

Every vertex now has at least three distinct neighbors and total degree four.
It can therefore meet at most one parallel pair.
\end{proof}

This reduction identifies a concrete remaining class: four-regular graphs whose doubled edges form a matching.
It does not justify contracting an arbitrary doubled edge.
That operation also changes cycles that use both endpoints without using the doubled edge.

\subsection{Equality is not unique}

The doubled $n$-cycle has $2^n+n$ cycles.
There is another equality example:
\[
 c(K_5)=\binom53+3\binom54+\frac{(5-1)!}{2}
       =10+15+12=37=2^5+5.
\]
Thus an exact equality classification cannot contain doubled cycles alone.
Proposition~\ref{prop:double-expansion} gives a limited classification.
If an equality graph has a vertex with two distinct neighbors, it must be a doubled cycle.
No asymptotic stability classification follows from this observation.

\subsection{Finite verification}

The script \texttt{check\_four\_regular.py} fixes one Hamilton cycle and enumerates the second up to reversal.
Every Hamilton-decomposable graph occurs, possibly more than once.
Cycle counts use vertex-simple cyclic orders, weighted by edge multiplicities.
For orders through eight, an independent implementation enumerates the binary cycle space and tests connectedness and degrees.
The two counts agree for every tested graph.
The script also checks $2{,}050$ doubled-path expansions and the equality criterion in Proposition~\ref{prop:double-expansion}.

\begin{center}
\begin{tabular}{r|rrrrrrr}
$n$&3&4&5&6&7&8&9\\\hline
Second Hamilton cycles&1&3&12&60&360&2520&20160\\
Maximum cycle count&11&20&37&70&135&264&521
\end{tabular}
\end{center}

These maxima equal $2^n+n$.
Together with complementation, the two-vertex base case, and the component argument, this verifies the bound through order nine.
The enumeration covers $23{,}116$ second-cycle orderings.
It is reproducible evidence, not a Lean certificate or an infinite-order proof.
The data are in \texttt{four-regular-results.json}.

An additional experiment tested $52{,}145$ contractions of arbitrary doubled edges.
The inequality $c(G)\le c(G/\{u,v\})+2^{n-1}+1$ held in every test.
Here the contraction identifies $u,v$ and deletes the two resulting loops.
This inequality remains unproved and is not used above.
